\documentclass[12pt, a4paper]{article}

% --- Pakiety ---
\usepackage[utf8]{inputenc}
\usepackage[T1]{fontenc}
\usepackage[polish]{babel}
\usepackage{geometry}
\usepackage{amsmath}
\usepackage{amssymb}
\usepackage{graphicx}
\usepackage{booktabs}
\usepackage{siunitx}
\usepackage{caption}
\usepackage{subcaption}
\usepackage{float}
\usepackage{hyperref}
\usepackage{xcolor}
\usepackage{array}
\usepackage{multirow}

% --- Marginesy ---
\geometry{
    top=2.5cm,
    bottom=2.5cm,
    left=2.5cm,
    right=2.5cm
}

% --- Konfiguracja siunitx ---
\sisetup{
    output-decimal-marker = {,},
    separate-uncertainty = true
}

% --- Dane raportu ---
\newcommand{\tytul}{Tytuł ćwiczenia}
\newcommand{\nrCwiczenia}{XX}
\newcommand{\autor}{Imię Nazwisko}
\newcommand{\nrGrupy}{X}
\newcommand{\dataWykonania}{\today}

% --- Tytuł dokumentu ---
\title{
    \large I Pracownia Fizyczna \\[0.5em]
    \textbf{Ćwiczenie nr \nrCwiczenia} \\[0.3em]
    \LARGE \tytul
}
\author{
    \autor \\
     Grupa \nrGrupy
}
\date{
    Data wykonania pomiarów: \dataWykonania \\
}

% =============================================
\begin{document}
% =============================================

\maketitle
\thispagestyle{empty}
\newpage

% -------------------------------------------
\section{Cel ćwiczenia}
% -------------------------------------------

Krótki opis celu eksperymentu — co mierzono i w jakim celu.

% -------------------------------------------
\section{Wstęp teoretyczny}
% -------------------------------------------

Opis teorii niezbędnej do zrozumienia ćwiczenia. Wzory kluczowe dla obliczeń, np.:

\begin{equation}
    F = ma
    \label{eq:newton}
\end{equation}

gdzie $F$ jest siłą wypadkową [\si{\newton}], $m$ — masą ciała [\si{\kilogram}], a $a$ — przyspieszeniem [\si{\meter\per\second\squared}].

% -------------------------------------------
\section{Opis układu pomiarowego}
% -------------------------------------------

Opis użytego sprzętu i schematu układu.

\begin{figure}[H]
    \centering
    % \includegraphics[width=0.6\textwidth]{schemat.png}
    \fbox{\parbox[c][5cm][c]{0.6\textwidth}{\centering Schemat układu pomiarowego}}
    \caption{Schemat układu pomiarowego.}
    \label{fig:schemat}
\end{figure}

Wykaz przyrządów:
\begin{itemize}
    \item Przyrząd 1 (dokładność: \SI{\pm 0.01}{\meter})
    \item Przyrząd 2 (dokładność: \SI{\pm 0.001}{\kilogram})
\end{itemize}

% -------------------------------------------
\section{Wyniki pomiarów}
% -------------------------------------------

\begin{table}[H]
    \centering
    \caption{Wyniki pomiarów.}
    \label{tab:pomiary}
    \begin{tabular}{c S[table-format=2.2] S[table-format=1.4] S[table-format=1.4]}
        \toprule
        \textbf{Nr} & \textbf{Wielkość 1 [\si{\meter}]} & \textbf{Wielkość 2 [\si{\second}]} & \textbf{Wynik [\si{\meter\per\second}]} \\
        \midrule
        1  &  &  &  \\
        2  &  &  &  \\
        3  &  &  &  \\
        4  &  &  &  \\
        5  &  &  &  \\
        \midrule
        Średnia & & & \\
        \bottomrule
    \end{tabular}
\end{table}

% -------------------------------------------
\section{Opracowanie wyników}
% -------------------------------------------

\subsection{Obliczenia}

Opis metody obliczeniowej. Przykładowe obliczenie wartości średniej:

\begin{equation}
    \bar{x} = \frac{1}{n} \sum_{i=1}^{n} x_i
\end{equation}

\subsection{Niepewności pomiarowe}

Niepewność standardowa typu A (statystyczna):

\begin{equation}
    u_A(x) = \sqrt{\frac{\sum_{i=1}^{n}(x_i - \bar{x})^2}{n(n-1)}}
\end{equation}

Niepewność standardowa typu B (z klasy przyrządu / podziałki):

\begin{equation}
    u_B(x) = \frac{\Delta x}{\sqrt{3}}
\end{equation}

Niepewność złożona:

\begin{equation}
    u_c(x) = \sqrt{u_A^2(x) + u_B^2(x)}
\end{equation}

Niepewność rozszerzona ($k = 2$, poziom ufności $\approx 95\%$):

\begin{equation}
    U(x) = k \cdot u_c(x) = 2 \cdot u_c(x)
\end{equation}

Propagacja niepewności dla wyniku $y = f(x_1, x_2, \ldots)$:

\begin{equation}
    u_c(y) = \sqrt{\sum_{i} \left(\frac{\partial f}{\partial x_i}\right)^2 u_c^2(x_i)}
\end{equation}

\subsection{Wynik końcowy}

\begin{equation}
    y = \bar{y} \pm U(y) = (\num{0.00} \pm \num{0.00}) \; \si{\meter\per\second}
\end{equation}

% -------------------------------------------
\section{Wykresy}
% -------------------------------------------

\begin{figure}[H]
    \centering
    % \includegraphics[width=0.8\textwidth]{wykres.pdf}
    \fbox{\parbox[c][7cm][c]{0.8\textwidth}{\centering Wykres}}
    \caption{Opis wykresu.}
    \label{fig:wykres}
\end{figure}

% -------------------------------------------
\section{Wnioski}
% -------------------------------------------

\begin{itemize}
    \item Podsumowanie uzyskanych wyników.
    \item Porównanie z wartością tablicową (jeśli dotyczy): $y_\text{tab} = \SI{0.00}{\meter\per\second}$, odchylenie: $X$\,\%.
    \item Dyskusja źródeł błędów i ocena jakości eksperymentu.
\end{itemize}

% =============================================
\end{document}
% =============================================
